\documentclass[11pt]{article}
\usepackage[T1]{fontenc}
\usepackage{amsmath,amssymb,amsthm}
\usepackage[margin=1in]{geometry}
\usepackage{microtype}
\usepackage[hidelinks]{hyperref}
\hypersetup{
  pdftitle={Quartic convergence in a diagonal counterexample to a claimed sixth-order inverse-matrix iteration},
  pdfauthor={Xiaojian Zeng},
  pdfsubject={Exact convergence-order counterexample for a norm-accelerated inverse-matrix iteration},
  pdfkeywords={matrix inversion, hyperpower iteration, Q-order, counterexample, exact verification}
}

\newtheorem{theorem}{Theorem}
\newtheorem{lemma}[theorem]{Lemma}
\newcommand{\C}{\mathbb C}
\newcommand{\norm}[1]{\lVert #1\rVert}
\DeclareMathOperator{\diag}{diag}

\title{Quartic convergence in a diagonal counterexample to a claimed\\
sixth-order inverse-matrix iteration}
\author{Xiaojian Zeng\\
PBC School of Finance, Tsinghua University, 100190\\
\texttt{xiaojianzeng@foxmail.com}}
\date{}

\begin{document}
\maketitle

\begin{abstract}
Cordero and Torregrosa recently proposed a norm-accelerated inverse-matrix
iteration and asserted order six for general nonsingular complex matrices.  We
give an exact counterexample under the update and residual convention of that
assertion.  For $A=I_3$, the Frobenius norm, and
$X_0=I_3-t\diag(1,2,3)$, $0<t\leq 10^{-2}$, the first residual contains a
nonzero quartic term.  A rational invariant-region argument then shows that
every fixed orbit in this interval has Q-order exactly four, with explicit
two-sided quartic bounds and a positive asymptotic error constant.  Thus the
quartic term is persistent rather than a one-step artifact.  We also identify
two invalid implications in the published order argument.  This note refutes
the stated general order-six conclusion; it neither asserts failure of
convergence nor classifies all orbits of the iteration.
\end{abstract}

\section{Context and residual map}

The Newton--Schulz iteration is a classical multiplication-based method for
matrix inversion \cite{Schulz1933}.  Its higher-order hyperpower relatives use
polynomials in the residual, and factorizations of those polynomials are used
to reduce the number of matrix products \cite{PetkovicPetkovic2015}.  Cordero
and Torregrosa \cite{CorderoTorregrosa2026} proposed a different scalar
norm-type acceleration, denoted CTM in their paper, and claimed order six for
every nonsingular complex matrix under a small-residual hypothesis.  The
purpose of this note is to audit that order assertion exactly.

Let $A\in\C^{n\times n}$ be nonsingular and let
\[
 R=I-AX,\qquad Y=X(2I-AX),\qquad
 \nu=\frac{\norm{I-AY}^{2}}{\norm{I-AX}^{2}},
\]
where $R\ne0$.  The CTM update is
\[
 X^+=Y\bigl((2+\nu)I-(1+\nu)AY\bigr).
\]
Since $I-AY=R^2$, direct expansion, without changing the residual convention,
gives
\begin{equation}
 \nu=\frac{\norm{R^2}^{2}}{\norm R^{2}},\qquad
 R^+=I-AX^+=-\nu R^2+(1+\nu)R^4.
\end{equation}
At $R=0$ the inverse has already been reached and the iteration terminates.
For comparison, the ordinary order-four hyperpower step has residual map
$R^+=R^4$.

We use the following quotient criterion for exact order, a standard sufficient
case within the Q-order framework \cite{Jay2001}: a nonterminating orbit is
said here to have Q-order exactly $p$ when
\[
 0<\lim_{k\to\infty}
 \frac{\norm{R_{k+1}}}{\norm{R_k}^{p}}<\infty.
\]
The norm below is always the Frobenius norm.

\section{The counterexample and its exact order}

\begin{theorem}\label{thm:main}
Apply \emph{(1)} with $A=I_3$ and
\[
 X_0=I_3-t\diag(1,2,3),\qquad 0<t\leq10^{-2}.
\]
Then
\begin{equation}
 R_1=t^4\diag(-6,-12,18)+7t^6\diag(1,16,81).
\end{equation}
The orbit is nonterminating and, for every $k\geq1$,
\begin{equation}
 \frac{27}{512}\norm{R_k}^{4}
 <\norm{R_{k+1}}<\frac43\norm{R_k}^{4}.
\end{equation}
Moreover, there are numbers
$a_\infty\in[3/10,2/5]$ and
$b_\infty\in[3/5,3/4]$ with
$a_\infty+b_\infty=1$ such that, on writing
\[
 m_\infty=\frac{a_\infty^4+b_\infty^4+1}
                  {a_\infty^2+b_\infty^2+1},\qquad
 S_\infty=(1+a_\infty^2+b_\infty^2)^{1/2},
\]
one has
\begin{equation}
 \lim_{k\to\infty}\frac{\norm{R_{k+1}}}{\norm{R_k}^{4}}
 =\frac{1-m_\infty}{S_\infty^3}\in(0,\infty).
\end{equation}
Consequently every fixed orbit in the displayed initial interval has Q-order
exactly four.
\end{theorem}

\begin{proof}
Put $D=\diag(1,2,3)$.  For $R_0=tD$,
\[
 \nu_0=\frac{t^4(1+16+81)}{t^2(1+4+9)}=7t^2.
\]
Coordinatewise substitution in (1) gives (2).  In particular,
\begin{equation}
 \frac{\norm{R_1}}{\norm{R_0}^{6}}
 =\frac{\norm{(-6+7t^2,-12+112t^2,18+567t^2)}}
        {14^3t^2}\longrightarrow\infty
\end{equation}
as $t\downarrow0$.  This already precludes a uniform local sixth-order
estimate.  We next prove that the quartic term persists on each fixed orbit.

After the first step write
\begin{equation}
 R_k=c_k\diag(-a_k,-b_k,1),\qquad q_k=c_k^2.
\end{equation}
The positivity and validity of this representation for all $k\geq1$ will
follow from the enclosure below.  For variables $(a,b,q)$ set
\[
 m=\frac{a^4+b^4+1}{a^2+b^2+1},\qquad d=1-m+qm.
\]
Equation (1), applied coordinatewise, induces the exact projective recurrence
\begin{equation}
 \begin{split}
 a^+&=\frac{a^2(m-a^2-qma^2)}d,\\
 b^+&=\frac{b^2(m-b^2-qmb^2)}d,\\
 q^+&=q^4d^2.
 \end{split}
\end{equation}
At $q=0$, call the first two components $P_1,P_2$.  Expansion gives
\begin{equation}
 P_1(a,b)+P_2(a,b)=1,\qquad
 P(a,b)=(a,b)\quad\hbox{if }a+b=1.
\end{equation}
Thus the leading homogeneous quartic map retracts the sign chart onto a line
of fixed projective directions.

We use the rational region
\begin{equation}
 \mathcal B=\left\{(a,b,q):
 \frac3{10}\leq a\leq\frac25,\quad
 \frac35\leq b\leq\frac34,\quad
 0\leq e:=1-a-b\leq\frac1{100},\quad
 0<q\leq10^{-12}\right\}.
\end{equation}
We record the estimates needed for a no-exit proof.  If $(a,b,q)\in\mathcal
B$ and $L=a^2-a^4+b^2-b^4$, rational endpoint bounds give
\begin{equation}
 L\geq\frac{153}{500},\qquad
 1+a^2+b^2\leq\frac{689}{400},\qquad
 d\geq\frac{L}{1+a^2+b^2}>\frac16.
\end{equation}
Both sign-preserving numerators in (7) are positive.  Indeed, $b^2>a^2$
implies $m-a^2>0$, while
\[
 m-b^2=\frac{1-b^2+a^2(a^2-b^2)}{1+a^2+b^2}
 \geq\frac{139}{689};
\]
the terms multiplied by $q\leq10^{-12}$ do not change either sign.

More precisely, exact cancellation in (7) gives
\begin{equation}
 1-a^+-b^+
 =\frac{q(a^4+b^4+1)^2}
 {a^2-a^4+b^2-b^4+q(a^4+b^4+1)}.
\end{equation}
As $a^4+b^4+1<27/20$, (7), (10), and (11) imply
\begin{equation}
 0\leq e^+\leq6q,\qquad q^+\leq2q^4.
\end{equation}
The factorizations
\begin{align*}
 P_1-a&=\frac{-a(a+1)(a-b^2)(a-b-1)(a+b-1)}
                   {a^4-a^2+b^4-b^2},\\
 P_2-b&=\frac{-b(b+1)(a^2-b)(a-b+1)(a+b-1)}
                   {a^4-a^2+b^4-b^2}
\end{align*}
and the bounds in (9)--(10) yield
\begin{equation}
 |P_1-a|\leq e,\qquad |P_2-b|\leq3e.
\end{equation}
Finally, differentiating (7) with respect to $q$ while fixing $a,b$ gives
\[
 \left|\frac{\partial a^+}{\partial q}\right|
 =\frac{a^2m^2(1-a^2)}{d^2}\leq6,\qquad
 \left|\frac{\partial b^+}{\partial q}\right|
 =\frac{b^2m^2(1-b^2)}{d^2}\leq21.
\]
Consequently
\begin{equation}
 |a^+-a|\leq e+6q,\qquad |b^+-b|\leq3e+21q.
\end{equation}

It remains to verify entry into the region and close the no-exit argument.
Put $u=t^2$.  Formula (2) gives
\begin{equation}
 a_1=\frac{6-7u}{18+567u},\quad
 b_1=\frac{12-112u}{18+567u},\quad
 e_1=\frac{686u}{18+567u},\quad
 q_1=u^4(18+567u)^2.
\end{equation}
For $0<u\leq10^{-4}$, direct cross-multiplication shows
\begin{equation}
 \frac{33}{100}\leq a_1\leq\frac13,\quad
 \frac{13}{20}\leq b_1\leq\frac23,\quad
 0\leq e_1\leq\frac1{250},\quad q_1<10^{-12}.
\end{equation}
While an orbit remains in $\mathcal B$, (12) gives
\[
 \sum_{j\geq1}q_j<2q_1,\qquad
 \sum_{j\geq1}e_j<\frac1{250}+12q_1.
\]
Together with (14), this bounds the total variations by
\begin{equation}
 \sum_{j\geq1}|a_{j+1}-a_j|<\frac1{200},\qquad
 \sum_{j\geq1}|b_{j+1}-b_j|<\frac{13}{1000}.
\end{equation}
The margins in (16) and (17) exclude a first exit through any of the four
ratio faces of (9); (11) and (12) exclude exits through the defect and scale
faces.  Hence the orbit never leaves $\mathcal B$.  In particular, (10) and
the numerator signs show inductively that (6) is valid and nonzero at every
step.

Let $S_k=(1+a_k^2+b_k^2)^{1/2}$.  From (7), (9), and (10),
\begin{equation}
 c_{k+1}=c_k^4d_k,\qquad \frac16<d_k<1,\qquad
 1<S_k<\frac43.
\end{equation}
Therefore
\[
 \frac{\norm{R_{k+1}}}{\norm{R_k}^4}
 =\frac{d_kS_{k+1}}{S_k^4},
\]
and (3) follows immediately.

The scale estimate in (12) gives $q_k\to0$, hence $R_k\to0$.  The summable
variations in (17) make $(a_k,b_k)$ Cauchy.  Equations (11) and (12) give
$1-a_\infty-b_\infty=0$, while (9) keeps both limiting coordinates strictly
between zero and one.  Thus
\[
 1-m_\infty
 =\frac{a_\infty^2-a_\infty^4+b_\infty^2-b_\infty^4}
        {1+a_\infty^2+b_\infty^2}>0.
\]
Taking the limit in (19) proves (4), and hence exact Q-order four.
\end{proof}

For every member of this family,
$\norm{R_0}=\sqrt{14}\,t<1$, so the initial data satisfy the stated hypothesis
of the source theorem.  They also satisfy the stronger restriction
$\norm{R_0}<3^{-1/3}$ used later in its convergence proof.  Because $A=I_3$,
the residual and inverse error have the same Frobenius norm.  Thus (5) refutes
the claimed general local order-six estimate under the source's own
convention, while Theorem~\ref{thm:main} determines the actual asymptotic
order on the explicit initial interval.

\section{The defect in the order-six argument}

Equation (2.7) of Cordero and Torregrosa \cite{CorderoTorregrosa2026} treats a
term shown only to be $O(\norm{e_k}^4)$ as
$\epsilon_k\norm{e_k}^4$ with $\epsilon_k\to0$.  A quartic big-O estimate
only makes the normalized coefficient bounded; it does not force that
coefficient to vanish.  Equations (2)--(4) show that it does not vanish on the
family above.

The following sentence makes a separate quantifier substitution.  From
``for every fixed $\epsilon>0$, eventually $\epsilon_k\leq\epsilon$,'' it sets
the moving value $\epsilon=\norm{e_k}^2$.  That substitution supplies no
single eventual index.  The desired conclusion would require the additional
estimate $\epsilon_k=O(\norm{e_k}^2)$, which is false on this family.

The result concerns the asserted order, not the convergence of CTM in general.
In particular, the displayed counterexample orbits converge rapidly, but with
Q-order four.  No claim is made here about other projective directions or a
classification of general CTM orbits.

\section*{Reproducibility and assistance disclosure}

All identities and inequalities used in the proof are exact.  A short
Python/SymPy verifier expands (2), (8), and (11), and checks the two derivative
identities used above; the rational bounds remain displayed in the proof and
do not rely on floating-point output.  This certificate-oriented presentation
is comparable in spirit to the exact low-dimensional matrix counterexample
and rational verification workflow in \cite{ZengLiuRatnavelu2026}.

OpenAI Codex assisted with proof exploration, exact symbolic checking,
adversarial mathematical review, literature discovery, and manuscript
preparation.  The author checked the stated argument and is responsible for
the correctness of the paper.  Computer algebra is used only to recompute
displayed algebraic identities, not as a substitute for the proof.

\bibliographystyle{plain}
\bibliography{references}

\end{document}
